%

\section{Continuous Action Control Using the MuJoCo Physics Simulator}
\label{sec:mujoco}

%
%
%
%
%
%
%
%
%
%
%
%
%
%
%
%
%
%
%
%
%
%
%
%
%
%
%
%
%
%
%
%
%
%
%
%
%
%
%
%
%

%
%
%
%
%
%
%
%
%
%
%
%
%
%
%

%
%
%
%
%
%
%
%
%
%

To apply the asynchronous advantage actor-critic algorithm to the Mujoco tasks the
necessary setup is nearly identical to that used in the discrete action domains, so
here we enumerate only the differences required for the continuous action domains.
The essential elements for many of the tasks (i.e. the physics models and task objectives)
are near identical to the tasks examined in \citep{lillicrap2015continuous}.
However, the rewards and thus performance are not comparable for most of the tasks due to
changes made by the developers of Mujoco which altered the contact model.

%
For all the domains we attempted to learn the task using the physical state as input.
The physical state consisted of the joint positions and velocities as well
as the target position if the task required a target.
In addition, for three of the tasks (pendulum, pointmass2D, and gripper) we also examined
training directly from RGB pixel inputs.
In the low dimensional physical state case, the inputs are mapped to a hidden state using one hidden layer with 200 ReLU units.
In the cases where we used pixels, the input was passed through two layers of spatial convolutions without any non-linearity or pooling.
In either case, the output of the encoder layers were fed to a single layer of 128 LSTM cells.
The most important difference in the architecture is in the the output layer of the policy
network. Unlike the discrete action domain where the action output is a Softmax, here the two outputs
of the policy network are two real number vectors which we treat
as the mean vector $\mu$ and scalar variance $\sigma^2$ of a multidimensional normal distribution
with a spherical covariance.
To act, the input is passed through the model to the output layer where we sample from the
normal distribution determined by $\mu$ and $\sigma^2$.
In practice, $\mu$ is modeled by a linear layer and $\sigma^2$ by a SoftPlus operation, $\log (1 + \exp (x))$,
as the activation computed as a function of the output of a linear layer.
In our experiments with continuous control problems the networks for policy network
and value network do not share any parameters, though this detail is unlikely to be crucial.
Finally, since the episodes were typically at most several hundred time steps long, we
did not use any bootstrapping in the policy or value function updates and batched each episode
into a single update.

As in the discrete action case, we included an entropy cost which encouraged exploration.
In the continuous case the we used a cost on the differential entropy of the normal distribution defined by
the output of the actor network, $-\frac{1}{2}(\log(2\pi\sigma^2)+1)$, we used a constant multiplier of $10^{-4}$ for this cost across all of the tasks
examined.
%
The asynchronous advantage actor-critic algorithm finds solutions for all the domains.  Figure \ref{fig-mujoco-wallclock}
shows learning curves against wall-clock time, and demonstrates that most of the domains from
states can be solved within a few hours. All of the experiments, including those done
from pixel based observations, were run on CPU.  Even in the case of solving the domains
directly from pixel inputs we found that it was possible to reliably discover solutions
within 24 hours.
Figure \ref{fig-mujoco-lr} shows scatter plots of the top scores against the sampled learning rates.
In most of the domains there is large range of learning rates that consistently achieve good performance
on the task.

%

%
%
%
%
%
%
%
%
%
%
